% Source-derived chapter 02: An overview of the methods
% Original source: rigid-local-systems-multiplicative-eigenvalue-problem
\chapter{An overview of the methods}
\label{rigid-local-systems-multiplicative-eigenvalue-problem-chapter-02}
\source{rigid-local-systems-multiplicative-eigenvalue-problem}

\begin{remark}[Source section]
\label{rigid-local-systems-multiplicative-eigenvalue-problem-chapter-02-source}
\source{rigid-local-systems-multiplicative-eigenvalue-problem}
\source{rigid-local-systems-multiplicative-eigenvalue-problem}
This chapter reproduces the corresponding section of the retrieved
LaTeX source, including its definitions, statements, examples, and
proofs. It has not yet been translated into Lean.
\notready
\end{remark}

% The source section title was: An overview of the methods
\label{uber}
\source{rigid-local-systems-multiplicative-eigenvalue-problem}
Using  some notation and definitions from the introduction, we outline the main constructions and the ideas that go into them.

The first input is the Mehta-Seshadri theorem \cite{MS} which recasts existence of unitary local systems on $\Bbb{P}^1-S$, in terms of non-emptiness of semi-stability loci, a condition involving existence of invariant sections of line bundes (in a GIT type situation). This is the same as existence of  non-zero sections of suitable tensor powers of a line bundle on a moduli stack of parabolic bundles (Theorem \ref{blacktea} is formulated for special unitary local systems, but the unitary case can be reduced to this). This condition is of the form
\begin{itemize}
\item
$H^0(\Par_{n,\mathcal{O},S},\mathcal{B}^m(\vec{\lambda},\ell))\neq 0$ for some $m$.
\end{itemize}

The next step is quantum saturation \cite{BHorn}, see  Proposition \ref{qTheorems} (1), which under  suitable conditions (of triviality of actions of isotropy groups), says that the above condition is equivalent to the non-vanishing for $m=1$.
We also have control over the rigidity of the moduli of unitary representations by quantum Fulton Proposition \ref{qTheorems}(2):  The moduli space is a point if and only if the above vector space is one dimensional for $m=1$ (note that the moduli space is the Proj of the section ring).

At this point, we have moved on from $P_n(s)$ the compact polyhedron controlling the unitary multiplicative eigenvalue problem to the cone over it, the effective cone of $\Par_{n,\mathcal{O},S}$ (see Theorem \ref{blacktea} and Proposition \ref{b9}). We now observe that some line bundles on $\Par_{n,\mathcal{O},S}$ satisfying geometrically defined properties give (some) extremal rays of the effective cones. These are the  F-line bundles defined in Definition \ref{defF}, see Lemma \ref{elgar}. The defining property is that they have a unique up to scalars non-zero section (even their tensor powers, but this follows from quantum Fulton), and further the zero scheme of this section is reduced and irreducible. The corresponding vertices of $P_n(s)$ are called F-vertices. Their defining property beyond being vertices of $P_n(s)$ is that the corresponding matrix equation has a unique solution up to  conjugation (See Definition \ref{defFprime}, and Lemma \ref{corresp}). Every line bundle $\mathcal{B}(\vec{\lambda},\ell)$ on $\Par_{n,\mathcal{O},S}$ has a level $\ell$ (Remark \ref{defL}).

The next input is strange duality. This says at the numerical level that $H^0(\Par_{n,\mathcal{O},S},\mathcal{B}(\vec{\lambda},\ell))$ has the same rank as a similar space of sections $H^0(\Par_{\ell,\widetilde{\mathcal{N}},S}\mathcal{B}(\vec{\lambda}^T,n))$ with $n$ and $\ell$ interchanged, and Young diagrams $\lambda_i$ replaced by their transposes, but the vector bundles (now of rank $\ell$)  now no longer have trivial determinants, but of a degree determined by the initial data, as in equation \eqref{GepnerPP}. We return to the history of \eqref{GepnerPP}, and the ideas behind understanding it geometrically in Section \ref{understanding}.


Once we have this, we have an equivalence of existence of a rank $\ell$ unitary local system with local monodromies $n$th roots of unity with that of a rank $n$ special unitary local system with corresponding local monodromies. The notions of F-line bundles and F-vertices are on the rank $n$-side. The notion that we consider on the rank $\ell$-side is Katz's notion of rigidity (the  representation should be irreducible, and the moduli space a point, also see Remark \ref{stretch2}).

The equivalence of the earlier paragraph leads to the question on what it does to rigid local systems on the rank $\ell$-side. This equivalence is non-linear since addition of Young diagrams as representations of $SL(n)$ is a row by row addition, and the transpose operation makes it a column by column operation. The answer is that rigids correspond to F-vertices, Theorem \ref{egregium}. %These are vertices of the polyhedron $P_n(s)$ for which the unitary matrix multiplication problem has a unique solution (but block diagonal).

We also associate to a rigid rank $\ell$ unitary representation with local monodromies $n$th roots of unity, a divisor on $\Par_{n,\mathcal{O},S }$ with enumerative significance (Lemma \ref{divisorE}, Proposition \ref{propP}), which puts non-trivial conditions (property (P) from the introduction) on the local monodromy of the rank $\ell$ unitary representation. One of the inputs here is a computation of divisor classes (see Section \ref{concept}).
The rigid unitary local system can then be constructed explicitly as a vector bundle of conformal blocks  equipped with the  Knizhnik-Zamalodchikov connection (Theorem \ref{KZe}).

We are therefore led to the question of determining vertices (and F-vertices)  of $P_n(s)$, which is the same as the question of determining the  extremal rays (and F-line bundles) of the cone $\Pic^+_{\Bbb{Q}}(\Par_{n,\mathcal{O},S})$ which is given by a system of inequalities parametrized by Gromov-Witten invariants  (Theorem \ref{ABW}, also see Proposition \ref{b9}). We determine the extremal rays on any regular face (see the discussion following Theorem \ref{ABW}) as follows:
\begin{enumerate}
\item[(a)] Some of extremal rays  correspond to explicit degeneracy loci  on $\Par_{n,\mathcal{O},S}$ (Theorems  \ref{Adagio} and \ref{brahms}). These give F-line bundles with numerical data given by Gromov-Witten invariants.
\item[(b)] The remaining extremal rays  on the face arise from a process of induction arising from a rational map
$\Par_{n,\mathcal{O},S}\to \ParL$ as explained in the introduction (Section \ref{schumann}). This process leads to a bijective parametrization of F-line bundles on the face  ( see Proposition \ref{condensed}).
\end{enumerate}
 The basic divisors in (a) correspond to the codimension 1 components in the na\"ive base locus of the rational map $\Par_{n,\mathcal{O},S}\to \ParL$ (see Remark \ref{explainmore}).
  \subsection{The strange duality construction}\label{understanding}
\source{rigid-local-systems-multiplicative-eigenvalue-problem}
The equality  \eqref{GepnerPP} is the numerical consequence of a canonical  duality of vector spaces:
\begin{equation}\label{canonD}
\source{rigid-local-systems-multiplicative-eigenvalue-problem}
H^0(\Par_{n,\mathcal{O},S},\ml)\leto{\sim} H^0(\Par_{\ell,\widetilde{\mathcal{N}},S},\ml')^*
\end{equation}
 This is the genus $0$ case of the strange duality conjecture for curves (\cite{Betrange,MO,Ou} and the surveys  \cite{paulyB, popa}). The strange duality conjecture has origins in mathematical physics (see e.g., \cite{NT}) where it is also called level-rank duality (also see \cite{BJDG}).

There is a natural divisor $D$ (called a theta divisor)  in the product $\Par_{n,\mathcal{O},S}\times \Par_{\ell,\widetilde{\mathcal{N}},S}$ whose associated line bundle is $\ml\boxtimes\ml'$. The divisor is a natural degeneracy locus: A point $(\mw,\me,\gamma)\times (\mt,\mf,\delta)$ is in $D$ if and only if there is a non-zero map $\phi:\mw\to \mt^*$ (or $\phi\in H^0(\Bbb{P}^1,\mt^*\tensor \mw^*)$) such that  $\phi_{p_i}(E^{p_i}_a)\subseteq G^{p_i}(\ell-\lambda^i_a)$ for $a=1,\dots,n$ and $i\in [s]$. Here $\mg\in \Fl_S(\mt^*)$ is induced from the identification  $\Fl_S(\mt^*)= \Fl_S(\mt)$ and the given $\mf\in\Fl_S(\mt)$. The expected dimension of such $\phi$ is zero in the settings considered in strange duality, so that the locus where a non-zero section exists is a divisor (see Section \ref{RamDiv} for one such setting).

The strange duality (or level-rank duality) is the assertion that $D$ results in an isomorphism
\eqref{canonD}. If we fix a point $x\in\Par_{\ell,\widetilde{\mathcal{N}},S}$, we get a corresponding section of $H^0(\Par_{n,\mathcal{O},S},\ml)$. Given the numerical equality of ranks, the strange duality is the assertion that  these sections as $x$ varies span $H^0(\Par_{n,\mathcal{O},S},\ml)$.


The theta divisor construction has been used frequently in the moduli theory of vector bundles (see e.g., \cite{Faltings}). In the non-parabolic (e.g., in higher genus) settings one is in situations where the Euler characteristic $\chi(\mt^*\tensor \mw^*)=0$, so that the locus where $H^0$ is non-zero is a divisor.

We  use only the numerical consequence of this duality. However a full picture of this duality may help in understanding the origin of the numerical equality (which is the first step in the proof of the vector space duality). The algebro-geometric proofs of strange duality proceed first by showing that the representation
theoretic spaces $H^0(\Par_{n,\mathcal{O},S},\ml)$ have ranks that are equal to enumerative numbers for the Grassmannian $\Gr(n,n+\ell)$ (because both are governed by ring structures,  the verification of this
comes down to very few cases, once the parameters are matched). Each point in the enumerative count is made to  give a section of $H^0(\Par_{n,\mathcal{O},S},\ml)$. The transversality in the enumerative count implies that these sections are linearly independent. Together with the numerical equality of the enumerative count and the rank of  $H^0(\Par_{n,\mathcal{O},S},\ml)$, we see that these sections span, and this gives the duality. A particularly simple case of this construction can be seen in the comparison of the structure constants in the representation rings of special linear groups and the intersection theory of Grassmannians \cite{BIMRN} (also see \cite[Section 6.2]{FBull}, and Remark \ref{classicall}).


\subsection{Enumerative computations}\label{concept}
\source{rigid-local-systems-multiplicative-eigenvalue-problem}
In Section \ref{enume} we have a computation of cycle classes of some natural loci in $\Par_{n,\mathcal{N},S}$.  Consider  the space formed by  the closure of $(\mw,\me,\gamma)$ such that the rank $n$ vector bundle $\mw$ has a subbundle $\mv$ of rank $r$ whose fibers in $\Gr(r,\mw_p)$ lie in given Schubert varieties $\Omega_{J^i}(\me_p)$ with respect to the given flags on $\mw_p$, with $p\in S$. We also assume that the expected dimension of such $\mv$ is $-1$ having fixed $(\mw,\me,\gamma)$. The set of  such $(\mw,\me,\gamma)$ gives  a natural divisor in $\Par_{n,\mathcal{N},S}$ (possibly zero).

In Proposition \ref{enumerative} we compute this divisor class, i.e., the corresponding line bundle $\mathcal{B}(\vec{\lambda},\ell)$. The coefficients that enter $\lambda_i$ and $\ell$ are obtained by weakening the number of conditions by one, so that the expected dimension of such $\mv$ is now zero (and we count over generic $(\mw,\me,\gamma)$). The weakening of conditions can change the degree of the underlying bundle since we allow for subbundles becoming subsheaves. The set of weakenings corresponds to enlargements of one of the Schubert varieties $\Omega_{J^i}(E_{\bull})$  for some $i$ by one dimension. Such enlargements correspond to $\Omega_{I}(E_{\bull})$ with $I$ obtained by replacing some $a\in J^i$ by $a+1$ assuming $a+1\notin J^i$ and $a\neq n$ (if $a=n$ we need a degree shift to handle this). The complexity of this last condition translates to property (P) in relation to the natural divisor on $\Par_{n,\mathcal{O},S}$ associated to a unitary rank $\ell$ rigid local system (as earlier in this section, Lemma \ref{divisorE} and Proposition \ref{propP}).




\iffalse
Let $\vec{J}=(J^1,\dots,J^s)$ be an $s$-tuple of subsets  of $[n]$ each of cardinality $r$. Recall from Definition \ref{cyclist} that the effective cycle $E=C(d,r,\mn,n,\vec{J})$ on $\Par_{n,\mn,S}$ is defined as a pushforward of
$\Omega(d,r,\mn,n,\vec{J})$ (see Definition \ref{mindful}).  We will compute $E$ when the codimension of this class is one, i.e.,
\begin{equation}\label{cdc}
\source{rigid-local-systems-multiplicative-eigenvalue-problem}
dn +\deg(\mn)r +r(n-r)-\sum_{i=1}^s |\sigma_{J^i}|=-1.
\end{equation}
We will use the following enumerative computation to prove Theorem \ref{brahms} in Section \ref{rumble}. It is also used to compare KZ local systems to strange duals of F-line bundles in Section \ref{EqKZ}. Recall the generalization of Gromov-Witten numbers given in  Definition \ref{GWgen}.
\begin{proposition}
\source{rigid-local-systems-multiplicative-eigenvalue-problem}
\notready\label{enumerative}
\source{rigid-local-systems-multiplicative-eigenvalue-problem}
Assume the codimension condition \eqref{cdc} and
write
$$\mathcal{O}(E)=\mathcal{B}(\vec{\lambda},\ell).$$
Write $\lambda^i=\sum_{b=1}^{n-1} c_{i}^{b}\omega_b$, where $\omega_b$ is the $b$th fundamental weight.
\fi
